\originalchapter{摄动与绝热不变量}
\originalsection{共振项与频率修正}
当方程含一个小参数, 直接令$x=x_0+\epsilon x_1+\cdots$通常只在有限时间内有效. 以$\ddot x+\omega_0^2x+\epsilon\alpha x^3=0$为例, 取$x_0=A\cos\omega_0t$. 因$\cos^3s=(3\cos s+\cos3s)/4$, 一阶方程出现与齐次解同频的强迫, 产生$t\sin\omega_0t$长期项. 这个项表示频率稍有偏差所积累的相位, 不能解释成真实振幅无限增长.

引入$\tau=\Omega t$, $\Omega=\omega_0+\epsilon\omega_1$, $x=A\cos\tau+\epsilon x_1(\tau)$. 一阶方程为
\[
\omega_0^2(x_1''+x_1)
=(2\omega_0\omega_1A-3\alpha A^3/4)\cos\tau
-\alpha A^3\cos3\tau/4.
\]
取$\omega_1=3\alpha A^2/(8\omega_0)$消去共振项, 非共振特解为$\alpha A^3\cos3\tau/(32\omega_0^2)$. 若$A$规定为初始最大位移, 还应加齐次项使$x_1(0)=0$, 即减去同系数的$\cos\tau$. 该近似要求$|\epsilon\alpha A^2|\ll\omega_0^2$; 不应在势失去稳定性或振幅很大时继续外推.

对摆$\ddot\theta+\omega_0^2(\theta-\theta^3/6+\cdots)=0$, 代入$\alpha=-\omega_0^2/6$得$\Omega=\omega_0(1-A^2/16+O(A^4))$, 所以振幅增大时周期增加. 这一符号可由能量积分中接近顶部运动变慢独立检查.

\originalsection{哈密顿平均}
设$H=H_0(I)+\epsilon H_1(I,\theta)$, $\theta$为$2\pi$周期角. 用生成函数$F_2=\theta I'+\epsilon S_1(\theta,I')$得到$K=H_0(I')+\epsilon[H_1+\omega(I')\partial_\theta S_1]+O(\epsilon^2)$. 若$\omega\ne0$, 可令
\[
\omega\partial_\theta S_1=\langle H_1\rangle-H_1,\qquad
\langle H_1\rangle=\frac1{2\pi}\int_0^{2\pi}H_1\dd\theta.
\]
右端均值为零, 积分产生周期$S_1$. 新作用量在一阶平均方程中守恒, 频率修正为$\epsilon\partial_I\langle H_1\rangle$. 多角系统中Fourier项的分母是$\vect k\cdot\vect\omega$; 接近共振时分母很小, 这种消去办法失效, 应保留慢共振角. 本册只在单频非共振范围内使用以上平均结果.

\originalsection{轨道进动}
自编例子取$U(r)=-\alpha/r+\epsilon\beta/r^2$. 力为$F=-\alpha/r^2+2\epsilon\beta/r^3$, Binet方程成为$u''+(1+2\mu\epsilon\beta/J^2)u=\mu\alpha/J^2$. 若括号正且摄动很小, 径向周期的角间隔为$\Delta\phi=2\pi/\sqrt{1+2\mu\epsilon\beta/J^2}$. 相对于一整圈的近日点进动为$\Delta\phi-2\pi\simeq-2\pi\mu\epsilon\beta/J^2$. $\beta>0$为反向进动, $\beta<0$为同向进动. 由于这里的摄动可以精确解出, 它也可检查一阶展开的系数和符号.

\originalsection{缓慢变化与绝热作用量}
若外部参数$\lambda(t)$变化时间远大于轨道周期, 每个冻结参数系统仍有周期轨道, 则$I=(2\pi)^{-1}\oint p\dd q$在平均意义下近似不变. 局部作用角变换随$\lambda$而动, 新哈密顿量含$\dot\lambda G(I,\theta,\lambda)$. 于是$\dot I=-\dot\lambda G_\theta$, 冻结$I,\lambda$作一周期平均时角导数均值为零. 这解释首阶漂移抵消的机制; 更长时间的误差控制需进一步假设平滑性、频率远离零及没有分离轨道穿越, 本册不把“变化慢”当无条件精确定理.

谐振子缓慢改变$\omega$时$I=E/\omega$近似固定, 故$E\propto\omega$, 而振幅$A=\sqrt{2I/(m\omega)}\propto\omega^{-1/2}$. 能量变化并不与守恒定律矛盾: 外部改变弹簧或约束需要作功. 当$\omega\to0$或摆穿越分离轨道, 周期不再比参数时间短, 绝热近似可失效.

本章来源: 8.09第5章. Duffing例子、摆的频率校验和可精确解的轨道摄动为自编推导, 未复制相关商业教材习题.
